\documentclass[11pt,a4paper]{article}
\usepackage[margin=2.3cm]{geometry}
\usepackage{amsmath,amssymb,bm}

\renewcommand{\v}[1]{\bm{#1}}
\newcommand{\ab}{\bar\alpha_s}
\newcommand{\dref}[1]{\texttt{\detokenize{#1}}}
\newcommand{\MSbar}{\overline{\rm MS}}
\title{Dilute--dilute limit, completed:\\ the soft region of the real emission with the plus-momentum cut,\\ and virtual $+$ real against the one-loop Lipatov vertex}
\author{verification note (scripts: \texttt{notes/scripts/dilute2/})}
\date{}
\begin{document}
\maketitle

\noindent\fbox{\parbox{0.97\textwidth}{\textbf{Verdict.}
\textbf{(1) Open point (c) of \texttt{ref3\_dilute} is closed:} the real emission of the paper (Groups II$+$III, dilute limit), integrated in $d_\perp=2-\epsilon$ with the cut $p^+>\Lambda$, supplies exactly $-\frac{\alpha_sN_c}{\pi}\big[(\frac2\epsilon-\ln\frac{\v k^2}{\mu^2})L+L^2\big]$, $L=\ln\frac{k^+}{\Lambda}$ (analytic, Eq.~(\ref{task1}); numerically the coefficient of the $\ln s$ behaviour is $2.0004$ against $2$, and $I(L)+L^2\to$ const). It cancels the extra term of Group~I. The double logarithm is the collinear logarithm $\ln(\v p_c^2)$, $\v p_c=\frac{p^+}{k^+}\v k$, of the soft gluon emitted next to the observed one, integrated over $p^+$ down to the cut; without the cut the same region gives the Fadin--Lipatov double pole.
\textbf{(2) Virtual $+$ real against the one-loop Lipatov vertex:} with the scheme conversion derived here (the CGC plus-momentum scheme is the BFKL scheme with energy scale $\v k^2$ on the projectile side and a fixed scale on the target side; non-factorisable term $\frac{\alpha_sN_c}{2\pi}\ln\v k^2\ln\frac{\v q_a^2}{\v q_b^2}$), the test ``$R_{\rm CGC}-R_{\rm FL}=F_P(\v q_a)+F_T(\v q_b)$'' was done at six points (three rectangles). \textbf{It is not passed as formulated:} the double differences are $-0.0174\pm0.0018$, $-0.014\pm0.005$, $+0.004\pm0.005$ (units of LO, $g=1$, $N_c=3$; $\frac{\alpha_sN_c}{\pi}=0.0760$), where a non-factorisable structure of size $0.10$ had to be reproduced: $83$--$87\%$ of it is there. The part that does not involve the scheme conversion (the structure of the one-loop vertex itself) is reproduced within errors ($-0.187\pm0.005$ against $-0.196$ and conversion $-0.005$).
\textbf{(3) The residual is not localised} to a row, and I cannot decide whether it is a discrepancy of the paper's finite terms or a missing piece of the target-side scheme conversion (the integrand of the rapidity logarithm on the target side depends on $\v k$ point by point, so the replacement of the cut $p^+>\Lambda$ by a rapidity cut is not guaranteed to be a $\v k$-independent redefinition of $\phi_T$; not computed). \textbf{The test of the finite terms against Fadin--Lipatov is therefore inconclusive: no agreement was established, and no error either.} Nothing was adjusted.}}

\section{Setting and notation}
This note continues \texttt{ref3\_dilute.tex} (same conventions): single inclusive gluon $(k^+,\v k)$ at mid rapidity, dilute projectile, soft modes $\Lambda<p^+<\vee$, pure glue, $d_\perp=2-\epsilon$, $\MSbar$ scale $\mu$, principal value in $\mathfrak B_3$; dilute target, $U=1+u$, first order in $u$ in the amplitude and in the conjugate amplitude; for the virtual part a point-like quark projectile. $\v q_a$ is the transverse momentum from the projectile, $\v q_b$ from the target, $\ab=\alpha_sN_c/\pi$, $V=\ln\frac{\vee}{k^+}$, $L=\ln\frac{k^+}{\Lambda}$, $2\omega(\v q)=\ab[\frac2\epsilon-\ln\frac{\v q^2}{\mu^2}]$,
\begin{align}
\frac{d^3N_{\rm LO}}{d^3k}=\mathcal N\int_{\v q_a}\frac{\phi_P(\v q_a)\phi_T(\v k-\v q_a)}{\v k^2}\,,\qquad
(\mathcal K\otimes\phi)(\v q)=\frac{\ab}{\pi}\int\frac{d^{2-\epsilon}p}{\v p^2}\,\phi(\v q+\v p)+2\omega(\v q)\phi(\v q)\,.
\end{align}
$\MSbar$ convention for transverse integrals (that of the scripts of \texttt{ref3\_dilute}): the measure is $d^{2-\epsilon}p\,(\pi e^{\gamma_E}\mu^2)^{\epsilon/2}$, so that $\frac1\pi\int_{|\v p|<\lambda}\frac{d^{2-\epsilon}p}{\v p^2}=-\frac2\epsilon+\ln\frac{\lambda^2}{\mu^2}$ ($\epsilon<0$ regulates the infrared). Inputs taken from \texttt{ref3\_dilute.tex}:
\begin{itemize}
\item[(R)] real emission (Groups II$+$III), exact identity with the quasi-multi-Regge matrix element $A_{gg}\equiv A_{\rm gluons}$ of BSV~\cite{BSV}:
\begin{align}
\frac{d^3N^{\rm(II+III)}}{d^3k}=\mathcal N\int_{\v q_a}\int_\Lambda^\vee\frac{dp^+}{p^+}\int d^{2-\epsilon}p\,\frac{\ab}{\pi}\,\phi_P(\v q_a)\phi_T(\v q_b)\frac{A_{gg}(\v q_a,\v q_b;\v k,\v p;x)}{\v q_a^2\v q_b^2}\,,\label{real}
\end{align}
with $\v q_b=\v k+\v p-\v q_a$, $x=\frac{k^+}{k^++p^+}$;
\item[(V)] Group I for a point-like quark, at $\mu=|\v k|$:
\begin{align}
\frac{\text{Group I}}{\rm LO}=2\omega(\v q_a)V+2\omega(\v q_b)L+\ab\Big[\frac2\epsilon L+L^2\Big]+C(\v q_a,\v q_b)\,.\label{virt}
\end{align}
By dimensional analysis (each pole $\frac2\epsilon$ comes with $-\ln\mu^2$, $C$ is dimensionless and the bare-coupling result has no other pole) the square bracket is $[(\frac2\epsilon-\ln\frac{\v k^2}{\mu^2})L+L^2]$ at general $\mu$, and $C$ is $\mu$ independent and invariant under a common rescaling and rotation of $(\v q_a,\v q_b)$.
\end{itemize}

\section{Task 1: the real emission at $|\v p|\lesssim\frac{\Lambda}{k^+}|\v k|$}\label{sec:soft}
\subsection{Analytic result}
Write $s\equiv p^+/k^+$ ($x=\frac1{1+s}$, $\frac{dp^+}{p^+}=\frac{ds}{s}$) and $\v p_c\equiv s\,\v k$, the collinear point ($x\v p_c=(1-x)\v k$, i.e.\ $\boldsymbol\Lambda=(1-x)\v k-x\v p=0$). Subtract the multi-Regge limit, which is the same function at both ends,
\begin{align}
\text{Eq.~(\ref{real})}&=(V+L)\,\mathcal N\!\int_{\v q_a}\!\int d^{2-\epsilon}p\,\frac{\ab}{\pi}\frac{\phi_P(\v q_a)\phi_T(\v q_b)}{\v k^2\v p^2}\notag\\
&\quad+\mathcal N\!\int_{\v q_a}\frac{\phi_P(\v q_a)\phi_T(\v k-\v q_a)}{\v k^2}\;\ab\int_{\Lambda/k^+}^{\vee/k^+}\frac{ds}{s}\,j(s)\,,\label{split}\\
j(s)&\equiv\frac1\pi\int d^{2-\epsilon}p\;\Big[\frac{\v k^2A_{gg}}{\v q_a^2\v q_b^2}-\frac{1}{\v p^2}\Big]\frac{\phi_T(\v q_b)}{\phi_T(\v k-\v q_a)}\,.\label{jdef}
\end{align}
The first term is the real part of $V\,\mathcal K\otimes\phi_P+L\,\mathcal K\otimes\phi_T$ (the two limits $x\to0,1$ are the same function of $\v p$; assigning it to the projectile or the target is a shift of the integration variable $\v q_a$). \texttt{ref3\_dilute} assumed that the $s$ integral of the second term converges at $s\to0$. It does not:

\paragraph{(i) Soft limit at fixed $s$.} For $s\to0$ the integral~(\ref{jdef}) is dominated by $|\v p|\sim|\v p_c|=s|\v k|\to0$: a soft gluon at a finite rapidity distance $\Delta y=y_p-y_k=\ln\frac{|\v p_c|}{|\v p|}$ from the observed one. There the matrix element is the eikonal factor for emission off the observed gluon and the two $t$-channel gluons (three adjoint charges, $T_i\cdot T_j=-\frac{N_c}2$),
\begin{align}
\frac{\v k^2A_{gg}}{\v q_a^2\v q_b^2}\ &\xrightarrow{\ s\to0,\ \v p=O(s)\ }\ \mathcal S\equiv\frac{1}{2\v p^2}\Big[1+\frac{\cosh\Delta y}{\cosh\Delta y-\cos\varphi}\Big]=\frac{\v p^2+\v p_c^2-\v p\cdot\v p_c}{\v p^2(\v p-\v p_c)^2}\,,\notag\\
\mathcal S-\frac{1}{\v p^2}&=\frac{\v p\cdot\v p_c}{\v p^2(\v p-\v p_c)^2}\label{softS}
\end{align}
($\varphi$ the angle between $\v p$ and $\v k$; corrections $O(s)$). The first term in the bracket is the emission between the two Reggeised gluons (the multi-Regge term), the second the two dipoles formed with the observed gluon; it is collinear singular at $\Delta y=\varphi=0$, i.e.\ at $\v p=\v p_c$, with unit residue. Eq.~(\ref{softS}) was checked numerically with the amplitude of the paper and with $A_{gg}$, also for $D\neq4$ (Sec.~\ref{sec:softnum}).
\paragraph{(ii) Transverse integral in $d_\perp=2-\epsilon$.} With $2\v p\cdot\v p_c=\v p^2+\v p_c^2-(\v p-\v p_c)^2$ and the vanishing of scaleless integrals,
\begin{align}
\frac1\pi\int d^{2-\epsilon}p\,\Big[\mathcal S-\frac1{\v p^2}\Big]&=\frac{\v p_c^2}{2\pi}\int\frac{d^{2-\epsilon}p}{\v p^2(\v p-\v p_c)^2}
=-\frac2\epsilon\,g(\epsilon)\Big(\frac{s^2\v k^2}{\mu^2}\Big)^{-\epsilon/2},\notag\\ g(\epsilon)&=\frac{e^{\gamma_E\epsilon/2}\,\Gamma(1+\frac\epsilon2)\Gamma(1-\frac\epsilon2)^2}{\Gamma(1-\epsilon)}=1-\frac{\pi^2}{48}\epsilon^2+\dots\label{softint}
\end{align}
so that, with $c(s)$ the angular average of the collinear residue,
\begin{align}
j(s)=-\frac2\epsilon\,c(s)+j_{\rm fin}(s)\,,\qquad c(s)\to1\,,\quad j_{\rm fin}(s)\to\ln\frac{s^2\v k^2}{\mu^2}\qquad(s\to0)\,.\label{jsmall}
\end{align}
The collinear pole of the real emission, taken at fixed $p^+$, has the scale $|\v p_c|=\frac{p^+}{k^+}|\v k|$, which runs with $p^+$: \emph{this} is the non-uniformity of the limit $\Lambda\to0$. For all $s$, $c(s)=[1-x(1-x)]^2\frac{\phi_T(\v k/x-\v q_a)}{\phi_T(\v k-\v q_a)}$ (found numerically to $10^{-8}$), i.e.\ $\frac{ds}{s}c(s)=dx\,\hat P(x)\,x^{-2}\,\frac{{\rm LO}(\v k/x)}{{\rm LO}(\v k)}$ with $\hat P(x)=\frac{1}{x(1-x)}-2+x(1-x)$: the DGLAP pole of the fragmentation of the parent gluon.
\paragraph{(iii) The $p^+$ integral with the cut.} From Eq.~(\ref{jsmall}), for $\vee\to\infty$ (the upper end converges, $j\sim s^{-0.8}$ numerically for the test function used),
\begin{align}
&\boxed{\;\ab\int_{\Lambda/k^+}^{\infty}\frac{ds}{s}\,j(s)=-\ab\Big[\Big(\frac2\epsilon-\ln\frac{\v k^2}{\mu^2}\Big)L+L^2\Big]\;-\;\frac2\epsilon\,\ab\,\mathcal P\;+\;\ab\,\mathcal R\;+\;O\big(L^2e^{-L}\big)\;}\label{task1}\\
&\mathcal P\equiv\int_0^\infty\frac{ds}{s}\,\big[c(s)-\theta(1-s)\big]\,,\qquad
\mathcal R\equiv\int_0^\infty\frac{ds}{s}\,\Big[j_{\rm fin}(s)-\theta(1-s)\ln\frac{s^2\v k^2}{\mu^2}\Big]\,,\notag
\end{align}
both $\Lambda$ independent. \textbf{The real emission of the paper, with its cut $p^+>\Lambda$, supplies exactly $-\ab[\frac2\epsilon L+L^2]$ (at $\mu=|\v k|$): the extra term of Group~I, Eq.~(\ref{virt}), cancels, pole, double logarithm and $\mu$ dependence.} The double logarithm is $\int_{e^{-L}}^1\frac{ds}{s}\ln s^2=-L^2$: the collinear logarithm $\ln\v p_c^2$ integrated over the plus momentum of the soft gluon; the pole term $-\frac2\epsilon L$ is the soft end, $\frac{2N_c}{1-x}$, of the DGLAP pole, cut at $1-x>\Lambda/k^+$. In the language of open point (c) of \texttt{ref3\_dilute}: for $|\v p|<\frac{\Lambda}{k^+}|\v k|$ the cut is at $\Delta y_{\min}=\ln\frac{\Lambda|\v k|}{k^+|\v p|}>0$ and removes the collinear point; the azimuthal average of Eq.~(\ref{softS}) is $\frac{1}{2\v p^2}[1+\coth|\Delta y|]$, and $\int_{\Delta y_{\min}}^\infty d\Delta y\,(\coth\Delta y-1)=-\ln(1-e^{-2\Delta y_{\min}})$ is finite there, while for $|\v p|>\frac{\Lambda}{k^+}|\v k|$ each $\v p$ carries the full collinear pole.

\paragraph{Relation to the calculation without cut (Fadin--Lipatov, BSV).} For $\epsilon<0$ the $s$ integral of Eq.~(\ref{softint}) converges at $s\to0$; the region removed by the cut is purely soft and
\begin{align}
\ab\int_{\Lambda/k^+}^{\infty}\frac{ds}{s}j(s)&=\ab\int_0^\infty\frac{ds}{s}j(s)-\ab\,\frac{2}{\epsilon^2}\,g(\epsilon)\Big(\frac{\v k^2}{\mu^2}\Big)^{-\epsilon/2}e^{\epsilon L}\notag\\
&=\ab\int_0^\infty\frac{ds}{s}j(s)-\ab\Big[\frac{2}{\epsilon^2}g\Big(\frac{\v k^2}{\mu^2}\Big)^{-\epsilon/2}+\Big(\frac2\epsilon-\ln\frac{\v k^2}{\mu^2}\Big)L+L^2\Big].\label{nocut}
\end{align}
The uncut real emission has the soft--collinear double pole $+\ab\frac{2}{\epsilon^2}(\v k^2/\mu^2)^{-\epsilon/2}$ $=\bar g_\mu^2\frac{2}{\epsilon_{\rm FL}^2}(\v k^2/\mu^2)^{\epsilon_{\rm FL}}$ ($\epsilon=-2\epsilon_{\rm FL}$, $\bar g^2_\mu\simeq\ab/4$), which cancels the $-\frac{2}{\epsilon^2}$ of the one-loop vertex; the plus-momentum cut trades it, in the real and in the virtual part separately, for $\frac2\epsilon L+L^2$. The two regularisations of the same integral differ by nothing else.

\subsection{Numerical verification}\label{sec:softnum}
Scripts \texttt{e1\_soft\_limit.py}, \texttt{e2\_real\_cut.py}, \texttt{e2b\_cut\_integral.py} (outputs \texttt{e1.out}, \texttt{e2\_m1.out}, \texttt{e2\_m2.out}, \texttt{e2b.out}); $\v k=(1,0)$, $\v q_a=(1.3,0.7)$, $\mu=|\v k|$.
\begin{itemize}
\item \emph{Soft limit, Eq.~(\ref{softS}).} Ratio $\v k^2A/(\v q_a^2\v q_b^2\,\mathcal S)$ at $\v p=s\,\v v$, $p^+=sk^+$, for the squared two-gluon amplitude of the paper at $O(u)$ (\texttt{d1.cgc\_sq}, Eq.~(\dref{eq:A2gamp})) and for $A_{gg}$ at $D=4,5,3.4$: $1-O(s)$ in all cases, e.g.\ for $\v v=(1.05,0.1)$ (next to the collinear point): $0.97876$, $0.99783$, $0.99978$, $0.99998$ at $s=10^{-2},\dots,10^{-5}$ (identical for the paper's amplitude and $A_{gg}$ to all digits; the $D$ dependence is $O(s^2)$).
\item \emph{Eq.~(\ref{softint}).} With a cut-off subtraction in two dimensions, $\frac1\pi\int d^2p\,[\mathcal S-\frac1{\v p^2}-\frac{\theta(|\v p-\v p_c|<\lambda)}{(\v p-\v p_c)^2}]=\ln\frac{\v p_c^2}{\lambda^2}$: $1.386294$ ($\lambda=0.5|\v p_c|$), $3.218876$ ($0.2|\v p_c|$), as predicted.
\item \emph{Full integrand at fixed $s$, Eq.~(\ref{jsmall}).} $j(s)$ of Eq.~(\ref{jdef}) with the exact $A_{gg}$, $\phi_T(\v q)=\v q^2e^{-\v q^2}$, $\v q_a$ fixed; the collinear singularity is subtracted with a Gaussian window, $\frac1\pi\int d^{2-\epsilon}p\,\frac{e^{-(\v p-\v p_c)^2/\rho^2}}{(\v p-\v p_c)^2}=-\frac2\epsilon+\ln\frac{\rho^2}{\mu^2}-\gamma_E$, the remainder integrated in two dimensions (three charts, Gauss--Legendre; changing the number of nodes by factors 2 and 3 changes $j_{\rm fin}$ by $4\times10^{-5}$ and $1\times10^{-5}$):
\begin{center}\small\setlength{\tabcolsep}{4pt}
\begin{tabular}{|c|cccccccc|}\hline
$\ln s$ & $-12$ & $-10$ & $-8$ & $-6$ & $-4$ & $-2$ & $0$ & $4$\\\hline
$c(s)$ & $0.999985$ & $0.999889$ & $0.999184$ & $0.99400$ & $0.95741$ & $0.76054$ & $0.63709$ & $0$\\
$j_{\rm fin}(s)$ & $-23.99985$ & $-19.99891$ & $-15.99335$ & $-11.9626$ & $-7.8185$ & $-3.5122$ & $-1.0667$ & $-0.0688$\\
$2\ln s$ & $-24$ & $-20$ & $-16$ & $-12$ & $-8$ & $-4$ & $0$ & \\\hline
\end{tabular}
\end{center}
A linear fit $j_{\rm fin}=a\ln s+b$ gives $a=2.0004$, $b=0.006$ on $-12\le\ln s\le-10$ ($a=2.0014$, $b=0.016$ on $[-12,-8]$; the deviation is the $O(s\ln^2s)$ correction): the coefficient of the pole is $c\to1$ and $j_{\rm fin}\to\ln s^2$, as claimed. (For $\ln s<-13$ the double-precision evaluation of $A_{gg}$ loses accuracy; not used.)
\item \emph{The cut integral, Eq.~(\ref{task1}).} $I(L)\equiv\int_{e^{-L}}^\infty\frac{ds}{s}j_{\rm fin}(s)$ (spline in $\ln s$ on the grid of \texttt{e2\_m2.out}):
\begin{center}
\begin{tabular}{|c|cccccc|}\hline
$L$ & $4$ & $5$ & $6$ & $8$ & $10$ & $12$\\\hline
$I(L)$ & $-16.2610$ & $-25.1331$ & $-36.0750$ & $-64.0391$ & $-100.0329$ & $-144.0317$\\
$I(L)+L^2$ & $-0.2610$ & $-0.1331$ & $-0.0750$ & $-0.0391$ & $-0.0329$ & $-0.0317$\\\hline
\end{tabular}
\end{center}
$I(L)+L^2$ converges to a constant ($\mathcal R\simeq-0.031$ for this $\phi_T$ and $\v q_a$; it is not universal and is not needed below). The coefficient of $L^2$ is $-1$ and there is no term linear in $L$ at $\mu=|\v k|$.
\end{itemize}

\section{Task 2: virtual $+$ real against the one-loop Lipatov vertex}
\subsection{The CGC result in the dilute limit}
Adding Eqs.~(\ref{virt}), (\ref{split}), (\ref{task1}) (point-like quark for the virtual part; $\vee/k^+\to\infty$, $\Lambda/k^+\to0$):
\begin{align}
\frac{d^3N}{d^3k}\Big|_{\rm dilute}=\mathcal N\int_{\v q_a}\Big\{&\frac{\phi_P(\v q_a)\phi_T(\v q_b)}{\v k^2}\Big[1+C(\v q_a,\v q_b)+\ab\Big(\mathcal R-\frac2\epsilon\mathcal P\Big)\Big]\notag\\
&+V\,\frac{(\mathcal K\otimes\phi_P)(\v q_a)\,\phi_T(\v q_b)}{\v k^2}+L\,\frac{\phi_P(\v q_a)\,(\mathcal K\otimes\phi_T)(\v q_b)}{\v k^2}\Big\},\label{cgc}
\end{align}
$\v q_b=\v k-\v q_a$. All rapidity logarithms multiply the full (infrared-finite) LO BFKL kernel; there is no $L^2$ and no pole other than the collinear one, $-\frac2\epsilon\ab\mathcal P$ (DGLAP; with the $\beta_0$ pole of the renormalised vertex it is the $\MSbar$ fragmentation counterterm). $\mathcal R$ is a functional of $\phi_T$ (non-local); $C$ multiplies $\phi_P\phi_T$ locally.

\subsection{The BFKL side and the rapidity schemes}
Fadin--Lipatov~\cite{FL}, in the form of BSV~\cite{BSV} (their Eqs.~\texttt{eq:news17nlo}, \texttt{MRKfinal}, \texttt{eq:kernelv}, \texttt{eq:defbsquare}): the production amplitude is
\begin{align}
\Gamma_A^{(s_{0;A})}\Big(\frac{s_{1}}{\sqrt{s_{0;A}\v k^2}}\Big)^{\omega(\v q_a)}\gamma^{(\v k^2,\v k^2)}(\v q_a,\v q_b)\Big(\frac{s_{2}}{\sqrt{\v k^2s_{0;B}}}\Big)^{\omega(\v q_b)}\Gamma_B^{(s_{0;B})},\label{regge}
\end{align}
$s_1=(p_{A'}+k)^2$, $s_2=(k+p_{B'})^2$, and the virtual kernel $\mathcal K^{(v)}_r$ (BSV \texttt{eq:kernelv}) is the interference of $\gamma^{(\v k^2,\v k^2)}$ with the Born vertex: the Reggeon scales are $s_{R;i,j}=\sqrt{s_{R;i}s_{R;j}}$ with $s_{R}=\v k^2$ for the produced gluon and $s_{0;A,B}$ for the external particles. I checked that the braces of \texttt{eq:kernelv} agree numerically, to $10^{-10}$ at random momenta, with twice the real part of the one-loop Lipatov vertex of Del Duca and Schmidt~\cite{DDS} (their Eq.~\texttt{looplip}) at Reggeisation scale $\tau=\v k^2$ (up to the pure constants $\propto\pi^2$, which depend on the normalisation factor $c_\Gamma$); so the scale convention and the printed finite terms are confirmed by an independent source:
\begin{align}
\frac{\mathcal K_r^{(v)}}{\mathcal K_r^{(B)}}&=\bar g_\mu^2\Big\{-\frac{2}{\epsilon_{\rm FL}^2}\Big(\frac{\v k^2}{\mu^2}\Big)^{\epsilon_{\rm FL}}+\frac{\beta_0}{N_c\,\epsilon_{\rm FL}}+X(\v q_a,\v q_b)\Big\},\notag\\
X&=\pi^2+\frac{3\v k^2}{\v q_a^2-\v q_b^2}\ln\frac{\v q_a^2}{\v q_b^2}-\ln^2\frac{\v q_a^2}{\v q_b^2}+\Big[\frac{\v k^2}{\v q_a^2-\v q_b^2}\Big(1-\frac{\v k^2(\v q_a^2+\v q_b^2-4\v q_a\!\cdot\!\v q_b)}{3(\v q_a^2-\v q_b^2)^2}\Big)\ln\frac{\v q_a^2}{\v q_b^2}\notag\\
&\qquad-\frac{\v k^2(\v q_a-\v q_b)^2}{6\v q_a^2\v q_b^2}+\frac{\v k^4(\v q_a^2+\v q_b^2)(\v q_a^2+\v q_b^2-4\v q_a\!\cdot\!\v q_b)}{6\v q_a^2\v q_b^2(\v q_a^2-\v q_b^2)^2}\Big]\qquad(n_f=0).\label{X}
\end{align}
For a massless projectile with plus momentum $P^+$ and a target with minus momentum $P^-$, $s_1=\frac{P^+}{k^+}\v k^2$, $s_2=2k^+P^-$. Regge factorisation requires $s_{0;A}$ ($s_{0;B}$) to be independent of the produced gluon; a change $s_{0;A}\to f(\v q_a)s_{0;A}$ is compensated in $\Gamma_A$. With $s_{0;A}=\v q_a^2$, $s_{0;B}=\v q_b^2$ the virtual logarithms of the cross section are
\begin{align}
2\omega(\v q_a)\Big[\ln\frac{P^+}{k^+}+\frac12\ln\frac{\v k^2}{\v q_a^2}\Big]+2\omega(\v q_b)\Big[\ln\frac{2k^+P^-}{\v k^2}+\frac12\ln\frac{\v k^2}{\v q_b^2}\Big],\label{FLlogs}
\end{align}
the rapidity differences of neighbouring particles. The real emission of the BFKL calculation covers the full phase space of the second gluon with the same matrix element $A_{gg}$ around the observed gluon (identity (R)); the sum of its MRK and QMRK pieces (BSV \texttt{eq:defbsquare} for $s_\Lambda\to\infty$, see Eq.~(\dref{d:schemeid}) of \texttt{ref3\_dilute}) is, at fixed $\v p$, $\frac{1}{\v k^2\v p^2}\times$(length of the rapidity interval) $+\int_0^\infty\frac{ds}{s}[\dots]$, i.e.\ the \emph{uncut} integral of Eq.~(\ref{nocut}), plus corrections at the ends of the interval which belong to the impact factors.

\paragraph{Conversion.} The CGC logarithms are plus-momentum logarithms on both sides: $V=\ln\frac\vee{k^+}$, $L=\ln\frac{k^+}{\Lambda}$, at fixed $\vee$, $\Lambda$, independent of $|\v k|$. Identify $\vee=P^+$ (a different constant is an impact-factor term) and define $Q_0^2\equiv2P^-\Lambda$ (the target is defined by all modes with $p^+<\Lambda$, i.e.\ $p^->\v p^2/(2\Lambda)$). Then
\begin{align}
V=\ln\frac{s_1}{\v k^2}\,,\qquad L=\ln\frac{s_2}{\v k^2}+\ln\frac{\v k^2}{Q_0^2}\,:
\end{align}
the CGC scheme is the BFKL scheme with $s_{0}=\v k^2$ on the projectile side and the fixed scale $s_0'=Q_0^2$ on the target side. Compared with Eq.~(\ref{FLlogs}):
\begin{itemize}
\item \emph{real part}: the rapidity interval of the second gluon is the same above the observed gluon; below it the CGC interval ends at $p^+=\Lambda$ instead of the kinematic limit, a difference $\ln\frac{\v p^2}{Q_0^2}$ in length, independent of $\v k$, which is by definition part of $\phi_T(\Lambda)$; in the soft region the cut gives Eq.~(\ref{nocut});
\item \emph{virtual part}: $[\text{CGC}]-[\text{Eq.~(\ref{FLlogs})}]=-\omega(\v q_a)\ln\frac{\v k^2}{\v q_a^2}+\omega(\v q_b)\ln\frac{\v k^2}{\v q_b^2}+2\omega(\v q_b)\ln\frac{\v q_b^2}{Q_0^2}$. Its pole part, $\frac{\ab}{\epsilon}\ln\frac{\v q_a^2}{\v q_b^2}$ plus a function of $\v q_b$, is of the form $F_P(\v q_a)+F_T(\v q_b)$ (it cancels against the $\v k$-independent real terms just described and the impact-factor terms). Its finite part is
\begin{align}
\frac{\ab}{2}\Big[\ln\frac{\v q_a^2}{\mu^2}\ln\frac{\v k^2}{\v q_a^2}-\ln\frac{\v q_b^2}{\mu^2}\ln\frac{\v k^2}{\v q_b^2}\Big]+F_T=\frac{\ab}{2}\ln\frac{\v k^2}{\mu^2}\,\ln\frac{\v q_a^2}{\v q_b^2}+F_P(\v q_a)+F_T(\v q_b)\,.\label{conv}
\end{align}
\end{itemize}
The first term of Eq.~(\ref{conv}) is \emph{not} of the form $F_P+F_T$: it is the price of evolving both sides in $p^+$ (the projectile in the natural, the target in the ``wrong'' variable), and has to be accounted for explicitly.

\paragraph{Prediction.} Subtracting, the soft--collinear double pole cancels between Eq.~(\ref{nocut}) and Eq.~(\ref{X}) with all its logarithms, $\mathcal R$ and $\mathcal P$ are common, and (for conventions see below)
\begin{align}
\boxed{\;D(\v q_a,\v q_b)\equiv C(\v q_a,\v q_b)-\frac{\ab}{4}X(\v q_a,\v q_b)+\frac{\ab}{2}\ln\frac{\v k^2}{\mu_0^2}\ln\frac{\v q_a^2}{\v q_b^2}+\frac{\ab\pi^2}{12}\;\overset{?}{=}\;F_P(\v q_a)+F_T(\v q_b)\;}\label{test}
\end{align}
($\v k=\v q_a+\v q_b$; $\mu_0$ any fixed unit). The functions $F_{P,T}$ are not determined by the present calculation: they depend on the definition of $\phi_{P,T}$ (impact factors, the identification of $\vee$, $Q_0$). The scheme-independent content of Eq.~(\ref{test}) is that the double difference over any rectangle in $(\v q_a,\v q_b)$ vanishes,
\begin{align}
{\rm DD}[D]\equiv D(\v q_{a},\v q_{b})+D(\v q_{a}',\v q_{b}')-D(\v q_{a},\v q_{b}')-D(\v q_{a}',\v q_{b})=0\,.\label{DD}
\end{align}
The constant $\frac{\ab\pi^2}{12}$ in Eq.~(\ref{test}) follows from $g(\epsilon)$, Eq.~(\ref{softint}), and the factor $\Gamma(1-\epsilon_{\rm FL})(4\pi)^{-\epsilon_{\rm FL}}$ in $\bar g_\mu^2$ (virtual double pole plus real soft--collinear term $=\frac{\ab\pi^2}{6}$ in the Fadin--Lipatov convention); it does not enter Eq.~(\ref{DD}). Likewise the $\beta_0$ pole: in the CGC Group~I with bare coupling the ultraviolet and collinear poles $\propto\frac{11}{3}$ cancel; after $\MSbar$ renormalisation of the coupling the pole is that of Eq.~(\ref{X}), with no finite remainder and no $\ln\mu^2$ in either expression.

\subsection{Numerical test}
Scripts \texttt{e0\_points.py}, \texttt{e3\_assemble.py}, \texttt{e3b\_prediction.py}, \texttt{e4\_localise.py} (outputs \texttt{e3.out}, \texttt{e3b.out}, \texttt{e4.out}; raw runs in \texttt{runs/}). The constant $C$ was computed with the unchanged scripts of \texttt{ref3\_dilute} (\texttt{v3c\_fast}, \texttt{v4b\_integrate}, \texttt{v5\_rowV}; same assembly as \texttt{v9b\_final}) at the corners of the grid $\v q_a\in\{(1.3,0.7),\,(0.5,0.4),\,(-0.4,0.9)\}$, $\v q_b\in\{(-0.3,-0.7),\,(0.5,-0.4)\}$ (labels $a_ib_j$; $a_1b_1$, $a_2b_2$ are the two points of \texttt{ref3\_dilute}); each corner is rotated and rescaled to $\v k=(1,0)$, $\mu=|\v k|$. Units: $g=1$, $N_c=3$, $\ab=0.075991$. The prediction for ${\rm DD}[C]$ (file \texttt{e3b.out}: $0.1143$, $-0.0770$, $-0.1913$) was committed before the new values of $C$ existed.

\paragraph{Checks at the new points.} At all six points the coefficients of the logarithms of Group~I are the trajectories: coefficient of $V$ against $2\omega(\v q_a)$: agreement to $10^{-6}$; coefficient of $L$ against $2\omega(\v q_b)$ (Row I.III fitted with the $L^2$ coefficient fixed to its analytic value): differences $(2,2,4,1,10,1)\times10^{-5}$; coefficient of $L^2$ (free fit) against $\ab$: differences $\le2\times10^{-5}$. This extends finding (b) of \texttt{ref3\_dilute} from two to six points.
\paragraph{Numerical accuracy of $C$.} The uncertainty quoted is the largest difference between the fit with the $L^2$ coefficient of Row~I.III fixed (central value), the free quadratic fit, the four three-point subsets, and the fit with the $L$ coefficient also fixed to $2\omega(\v q_b)$. Further checks at $a_1b_2$: Row~I.III at $\Lambda=10^{-5}$ with half the radial nodes and a different partition radius ($\rho=0.35$): $8.41873$ against $8.41919$ ($5\times10^{-4}$, \texttt{e6}); Row~I.V with 24 instead of 12 nodes and two more end-point values: $0.005919$ against $0.005924$; Row~I.IV: spline, monotone interpolation and Gauss weights agree to $3\times10^{-5}$. The point $a_3b_1$ is less accurate ($\pm0.005$: there the free fit of the $L^2$ coefficient is off by $2\times10^{-5}$).
\begin{center}
\renewcommand{\arraystretch}{1.15}
\begin{tabular}{|c|c|c|c|c|c|c|c|}\hline
point & $|\v k|$ & $C$ (CGC) & $\frac{\ab}{4}X$ (FL) & $C-\frac{\ab}4X$ & conv.\ term & $D-\frac{\ab\pi^2}{12}$ & $E$\\\hline
$a_1b_1$ & $1.0000$ & $0.0546\pm0.0006$ & $0.20986$ & $-0.1553$ & $0$ & $-0.1553$ & $-0.1612$\\
$a_1b_2$ & $1.8248$ & $0.1155\pm0.0015$ & $0.36408$ & $-0.2485$ & $+0.0764$ & $-0.1722$ & $-0.1686$\\
$a_2b_1$ & $0.3606$ & $0.1472\pm0.0005$ & $0.20258$ & $-0.0554$ & $+0.0269$ & $-0.0285$ & $-0.0380$\\
$a_2b_2$ & $1.0000$ & $0.3050\pm0.0008$ & $0.36783$ & $-0.0628$ & $0$ & $-0.0628$ & $-0.0628$\\
$a_3b_1$ & $0.7280$ & $0.1337\pm0.0048$ & $0.23039$ & $-0.0967$ & $-0.0124$ & $-0.1091$ & $-0.1035$\\
$a_3b_2$ & $0.5099$ & $0.1041\pm0.0004$ & $0.19959$ & $-0.0955$ & $-0.0441$ & $-0.1396$ & $-0.1245$\\\hline
\end{tabular}
\end{center}
(conv.\ term $=\frac{\ab}2\ln\v k^2\ln\frac{\v q_a^2}{\v q_b^2}$ in the common units of the grid; $E$ is defined in Eq.~(\ref{E}) below.)
\paragraph{Double differences, Eq.~(\ref{DD}).}
\begin{center}\small\setlength{\tabcolsep}{4pt}
\renewcommand{\arraystretch}{1.15}
\begin{tabular}{|c|c|c|c|c|c|}\hline
rectangle & ${\rm DD}[C]$ & ${\rm DD}[\frac{\ab}4X]$ & $-{\rm DD}[\text{conv}]$ & predicted ${\rm DD}[C]$ & ${\rm DD}[D]$\\\hline
$(a_1,a_2)\times(b_1,b_2)$ & $+0.0969\pm0.0018$ & $+0.0110$ & $+0.1033$ & $+0.1143$ & $-0.0174\pm0.0018$\\
$(a_1,a_3)\times(b_1,b_2)$ & $-0.0905\pm0.0051$ & $-0.1850$ & $+0.1080$ & $-0.0770$ & $-0.0135\pm0.0051$\\
$(a_2,a_3)\times(b_1,b_2)$ & $-0.1874\pm0.0049$ & $-0.1961$ & $+0.0048$ & $-0.1913$ & $+0.0039\pm0.0049$\\\hline
\end{tabular}
\end{center}
(errors added in quadrature; the linear sums are $0.0033$, $0.0073$, $0.0065$; the three rectangles are not independent, the third is the difference of the first two. With the free and with the fully constrained fits of Row~I.III the first entry of the last column is $-0.0178$ and $-0.0170$.)
\begin{itemize}
\item Without the conversion term the test would fail by $0.086$ and $0.095$ in the first two rectangles: the non-factorisable structure $\ln\v k^2\ln(\v q_a^2/\v q_b^2)$ required by the different rapidity variables \emph{is} present in the CGC result, with $83\%$ and $87\%$ of the coefficient derived in Eq.~(\ref{conv}).
\item In the third rectangle the conversion term is small and the non-factorisable structure of the one-loop vertex itself is large (${\rm DD}[\frac{\ab}{4}X]=-0.196$, mostly from $\frac{3\v k^2}{\v q_a^2-\v q_b^2}\ln\frac{\v q_a^2}{\v q_b^2}$ and the bracket of Eq.~(\ref{X})): the CGC constant reproduces it, ${\rm DD}[C]=-0.187\pm0.005$ against $-0.191$.
\item In the first rectangle, where the numbers are most accurate, the double difference is $-0.0174\pm0.0018$ ($0.23\,\ab$): \textbf{not zero}.
\end{itemize}
\paragraph{Stronger form of the test.} $C-\frac{\ab}4X$ is invariant under rotations and under a common rescaling of $(\v q_a,\v q_b)$. Then $D=F_P+F_T$ and the scaling behaviour of the conversion term force $F_P=\frac{\ab}4\ln^2\v q_a^2+c\ln\v q_a^2$, $F_T=-\frac{\ab}{4}\ln^2\v q_b^2-c\ln\v q_b^2$ up to constants, i.e.
\begin{align}
E\equiv C-\frac{\ab}{4}X+\frac{\ab}{4}\ln\frac{\v q_a^2}{\v q_b^2}\,\ln\frac{\v k^4}{\v q_a^2\v q_b^2}\overset{?}{=}c\,\ln\frac{\v q_a^2}{\v q_b^2}+e_0\label{E}
\end{align}
with two numbers $c$, $e_0$ for all kinematics (in particular $E$ must not depend on the angle between $\v q_a$ and $\v q_b$). A fit to the six points gives $c=-0.0675$, $e_0=-0.0645$ with residuals $(-7.3,\,8.7,\,3.1,\,1.7,\,-4.3,\,-1.9)\times10^{-3}$, rms $5.2\times10^{-3}$, against errors of $(0.6,1.5,0.5,0.8,4.8,0.4)\times10^{-3}$: Eq.~(\ref{E}) does not hold either.%%SYM%%

\subsection{Localisation of the residual (Task 3)}
\begin{enumerate}
\item \emph{Not the numerics of $C$} at the level of $0.017$ (accuracy checks above; the log coefficients are right at all points).
\item \emph{Not the transcription or scale convention of the one-loop vertex:} BSV's $\mathcal K^{(v)}_r$ equals the Del Duca--Schmidt vertex at $\tau=\v k^2$ to $10^{-10}$; the scale assignment $s_{R}=\sqrt{\v k^2s_0}$ is the one stated by BSV, and the $\v k$ dependence of the single poles it implies, $+\frac{\ab}{\epsilon}\ln\v k^2$ in the sum of double pole and trajectory terms, agrees with the universal infrared structure of one-loop amplitudes (colour charges $T_{q_a}\cdot T_g=T_{q_b}\cdot T_g=-\frac{N_c}2$) and with the poles of the CGC result after the conversion.
\item \emph{Not the soft regions of the loop} (scripts \texttt{e7*}, outputs \texttt{e7\_a1b1.out}, \texttt{e7b.out}, \texttt{e7cd.out}). A hard cut $s=p^+/k^+>\Lambda/k^+$ differs from dimensional regularisation also in the region $\v p\to0$ at fixed $u=\v p^2/(s\v k^2)=p^-/k^-$ (loop gluons with $p^+\lesssim\Lambda$ whose light-cone energy is comparable with that of the observed gluon; the scale $\sqrt{\Lambda/k^+}|\v k|$ of the numerics of Row~I.III), by a term $\frac{1}{4\pi}\int\frac{du}{u}\ln u\,[\dots]$ which Eq.~(\ref{nocut}) does not contain. I checked the integrands: in this region Rows~I.III and I.IV are independent of $u$ after the angular average ($-0.72443$ and $-\frac1{3\pi}$ at $a_1b_1$), and the resolved Row~I.V is, to $10^{-6}$, $\text{const}+8\pi\ab\,[\frac{P_k}{1+u}-\frac{P_q}{1+u\v k^2/\v q_a^2}]$, single-scale functions with the weights $P_k$, $-P_q$ of the two terms of the Lipatov vertex; their $u$ dependence and the $\ln^2\frac{\v q_a^2}{\v k^2}$ it would generate must be cancelled by the commutator part of Row~I.II, which is a single-scale function with the same weights (the net step has to vanish for the double pole to be the universal one, and the measured $L$ coefficient leaves no room for a term $P_q\ln\frac{\v q_a^2}{\v k^2}$; inferred, the integrand of Row~I.II is not available point by point). In the soft--collinear region $\v p=s\v v$ the sum of Rows~I.III, I.IV and the resolved I.V is the same function of $\v v$ at three kinematic points to $10^{-5}$ (the individual rows differ by factors). So in the soft regions the hard cut differs from dimensional regularisation by the universal term of Eq.~(\ref{nocut}) and constants only.
\item \emph{Character of the residual.} It sits in the part that is odd under $\v q_a\leftrightarrow\v q_b$, where the rapidity schemes enter, and not in the structure of the vertex (third rectangle). Adding one more structure to Eq.~(\ref{E}) (\texttt{e4.out}): no single candidate brings the rms residual below $1.7\times10^{-3}$ (best: $P_q\ln\frac{\v q_a^2}{\v q_b^2}$, $\ln^2\frac{\v k^2}{\v q_b^2}$); reducing the coefficient of the conversion term from $\frac12$ to $0.43$ on both sides gives $2.5\times10^{-3}$. No clean identification.
\item \emph{What could not be excluded.}
(a) \emph{The target-side conversion beyond the leading regions.} Eq.~(\ref{conv}) assumes that the loop gluons with $p^+<\Lambda$ (down to the kinematic limit of a physical target, $p^+>\v p^2/2P^-$) change the result by a $\v k$-independent functional of $\phi_T$. In the gauge and variables of the paper the integrand of the target-side logarithm at fixed loop momentum, $\tau(\v p;\v k,\v l)=\lim_{s\to0}s\,[F_{\rm III}+F_{\rm IV}+F_{\rm V}]$, depends on $\v k$ at fixed $\v l$ (\texttt{e7cd.out}: e.g.\ $-2.45$, $-9.82$, $-2.69$, $-3.01$ at $\v p=(0.4,0.5)$, $\v l=(0.3,0.7)$ for four $\v k$; only its integral is the trajectory). A weight $\ln\v p^2$ under this integral, which is what the difference between a cut in $p^+$ and a cut in rapidity or $p^-$ amounts to, is therefore not obviously of the form $F_T(\v q_b)$. This is the same physics as the ``kinematical constraint'' of plus-momentum-ordered evolution on the target side. I did not compute it (it needs the integrand of the commutator part of $\mathfrak A^{(3)}$, which the paper gives only integrated, and a definition of the cut for each row); it is a term of exactly the type of the residual.
(b) \emph{Row I.III ($\mathfrak B_3$, principal value).} Rows~I.IV and I.V are built from coefficients ($\mathfrak A$, $\mathfrak B_2$, $\mathfrak C_1$) which the real emission validated to $10^{-15}$; the kinematics-dependent part of the Row~I.II constant cancels against Row~I.V; $\mathfrak B_3$ is the only coefficient that enters $C$ without an independent check.
(c) \emph{The $\{t^a,t^b\}$ (non-planar, $1/N_c^2$) terms of the quark projectile}; the residual is $0.23\,\ab$, a repetition with $SU(2)$ would separate them (not done).
(d) Errors in the evaluation of $C$ in \texttt{ref3\_dilute} that do not affect the logarithms.
\end{enumerate}

\section{Verdict, limitations, open points}
\begin{itemize}
\item \textbf{Task 1: settled.} The real emission of the paper with the cut $p^+>\Lambda$ contains, besides $\ln\frac{\vee}{\Lambda}$ times the real BFKL kernel, the term $-\ab[(\frac2\epsilon-\ln\frac{\v k^2}{\mu^2})L+L^2]$, the DGLAP pole $-\frac2\epsilon\ab\mathcal P$ and a finite $\Lambda$-independent remainder, Eq.~(\ref{task1}). The statement of \texttt{ref3\_dilute} that ``for the real emission the plus-momentum cut-offs require no conversion'' holds at fixed $\v p\neq0$ only and must be replaced by Eq.~(\ref{nocut}). Virtual $+$ real in the dilute limit is Eq.~(\ref{cgc}): all rapidity logarithms multiply the full LO BFKL kernel, there is no $\ln^2\frac{k^+}{\Lambda}$.
\item \textbf{Task 2: done, result negative as formulated.} With matched regulators and the explicit conversion of the rapidity schemes the difference between the CGC result (Group~I constant of \texttt{ref3\_dilute}, six points) and the one-loop Lipatov vertex is \emph{not} of the form $F_P(\v q_a)+F_T(\v q_b)$: ${\rm DD}=-0.0174\pm0.0018$ on the best rectangle, to be compared with a non-factorisable structure of $0.103$ that is otherwise reproduced. The vertex-type structure is reproduced within $0.005$.
\item \textbf{Task 3: not localised.} Excluded: numerics, the Fadin--Lipatov input, scheme artefacts of the soft regions. Not excluded: the target-side conversion of the hard $p^+$ cut beyond the leading regions (item 5a above), $\mathfrak B_3$ (Row~I.III), the non-planar terms of the quark projectile, an error of the earlier evaluation. \textbf{The paper's finite terms have neither passed nor failed the test: it is inconclusive}, because the conversion of the target-side cut-off is established only as far as pole and logarithmic terms and the soft regions go; within that conversion there is a residual of $\approx15\%$ of the scheme term ($0.2\,\ab\times$LO).
\end{itemize}
\paragraph{Limitations.} (i) Virtual part for a point-like quark projectile only (the $f^{abc}\{\rho^b,\rho^c\}$ part of $\mathfrak A^{(3)}$ vanishes there). (ii) The check is at the level of the amplitudes and coefficients of the paper as coded in \texttt{ref3\_dilute}, not of the coordinate-space finite terms as printed. (iii) $C$ is obtained from fits in $L$ ($10^{-5}\ge\Lambda\ge10^{-8}$); a $\Lambda$-independent quadrature error of Row~I.III would not be seen in the fit (one grid check: $5\times10^{-4}$). (iv) The $\mu$ dependence used ($C$ independent of $\mu$, bracket of Eq.~(\ref{virt}) with $\ln\frac{\v k^2}{\mu^2}$) follows from dimensional analysis, it was not recomputed. (v) $F_P$, $F_T$ and the constant are not tested (they depend on the impact factors). (vi) $\mathcal R$ and $\mathcal P$ are common to both sides by the identity of the real integrands; they were not evaluated for the paper's printed real finite terms.
\paragraph{What would decide it.} (1) Recompute Group~I with the lower cut on the loop gluon imposed on $p^-$ (or rapidity) instead of $p^+$ in the region $p^+\ll k^+$, row by row (Rows~I.III, I.IV, I.V are available point by point; $\mathfrak A^{(3)}$ has to be redone before the $\xi$ integration), and repeat the test; this removes assumption 5(a). (2) Repeat at $N_c=2$ (separates the non-planar terms). (3) Three more points with $|\v q_a|=|\v q_b|$ at different angles: there every scheme term and $F_P+F_T$ reduce to a constant, so $C-\frac{\ab}{4}X$ must be one number; this tests the vertex part alone.

\section{Drop-in text for the paper}
Replacement for the last sentences of the ``Dilute limit'' paragraph proposed in \texttt{ref3\_dilute.tex} (from ``For a point-like quark projectile the virtual correction (Group I) is'' on), and for the corresponding open point. Macros of the draft.
\begin{verbatim}
For a point-like quark projectile the virtual correction (Group I) is
\begin{equation}\label{eq:dilute_virtual}
\frac{\text{Group I}}{\rm LO}=2\omega(\v q_a)\ln\frac{\vee}{k^+}
+2\omega(\v q_b)\ln\frac{k^+}{\Lambda}
+\frac{\alpha_sN_c}{\pi}\Big[\Big(\frac2\epsilon-\ln\frac{\v k^2}{\mu^2}
\Big)\ln\frac{k^+}{\Lambda}+\ln^2\frac{k^+}{\Lambda}\Big]+C(\v q_a,\v q_b),
\end{equation}
$2\omega(\v q)=\frac{\alpha_sN_c}{\pi}[\frac2\epsilon-\ln\frac{\v q^2}
{\mu^2}]$, with a finite function $C$ and no $\frac{67}{9}$ (checked
numerically at six kinematic points). The square bracket is cancelled by
the real emission: at fixed $p^+$ the transverse integral of
Eq.~\eqref{eq:dilute_real} has, besides the soft pole of the BFKL
kernel, a collinear pole at $\v p=\frac{p^+}{k^+}\v k$ whose scale runs
with $p^+$; for $p^+\ll k^+$ the integrand is the eikonal factor
$\frac{1}{\v p^2}+\frac{\v p\cdot\v p_c}{\v p^2(\v p-\v p_c)^2}$,
$\v p_c=\frac{p^+}{k^+}\v k$, and
\begin{equation}
\int_{\Lambda}\frac{dp^+}{p^+}\int\frac{d^{2-\epsilon}p}{\pi}
\Big[\frac{\v k^2A_{gg}}{\v q_a^2\v q_b^2}-\frac{1}{\v p^2}\Big]
=-\Big(\frac2\epsilon-\ln\frac{\v k^2}{\mu^2}\Big)\ln\frac{k^+}{\Lambda}
-\ln^2\frac{k^+}{\Lambda}+\dots,
\end{equation}
where the dots are the collinear (DGLAP) pole of the fragmentation of
the parent gluon and a $\Lambda$-independent remainder. Hence in the
sum of virtual and real contributions all logarithms of $\vee$ and
$\Lambda$ multiply the full LO BFKL kernel acting on $\phi_P$ and
$\phi_T$, and there is no $\ln^2\frac{k^+}{\Lambda}$; without the cut
the same region gives the soft--collinear double pole
$\epsilon^{-2}(\v k^2/\mu^2)^{\epsilon}$ of the real and virtual parts of
Ref.~\cite{Fadin:1998py}. The plus-momentum cut-offs correspond to the
BFKL energy scale $s_0=\v k^2$ on the projectile side and to a fixed
scale on the target side; relative to the symmetric scales of
Ref.~\cite{Fadin:1998py} this moves a term
$\frac{\alpha_sN_c}{2\pi}\ln\v k^2\ln\frac{\v q_a^2}{\v q_b^2}$ into the
finite part. With this term, $C$ reproduces the non-factorisable
structure of the one-loop central-emission
vertex~\cite{Fadin:1998py,DelDuca:1998cx,Bartels:2006hg} up to a
remainder of about $0.2\,\frac{\alpha_sN_c}{\pi}$ which is not of the
form $F_P(\v q_a)+F_T(\v q_b)$; whether this remainder is removed by the
exact conversion of the target-side cut-off (a cut in $p^+$ instead of
rapidity) has not been established.
\end{verbatim}
\noindent The sentence beginning ``The plus-momentum cut-offs correspond\dots'' and the following one report the outcome of the present note; the authors may prefer to state only that the comparison of the finite function $C$ with the one-loop vertex is not complete. In the list of open points of the paper, the item on the dilute limit should read:
\begin{verbatim}
Dilute limit. Real emission: exact agreement with the quasi-multi-Regge
matrix element; the cut p^+ > Lambda produces
-(alpha_s N_c/pi)[(2/eps - ln(k^2/mu^2)) L + L^2], which cancels the
corresponding term of Group I (point-like projectile). Open: the finite
function C of Group I agrees with the one-loop Lipatov vertex only up to
a non-factorisable remainder of about 0.2 alpha_s N_c/pi (15 per cent of
the rapidity-scheme term); the conversion of the target-side cut-off
beyond logarithmic accuracy, the N_c dependence, and Row I.III remain to
be checked.
\end{verbatim}

\begin{thebibliography}{9}
\bibitem{BSV} J.~Bartels, A.~Sabio Vera, F.~Schwennsen, JHEP 0611:051 (2006), hep-ph/0608154.
\bibitem{FL} V.S.~Fadin, L.N.~Lipatov, Phys.\ Lett.\ B429 (1998) 127, hep-ph/9802290.
\bibitem{DDS} V.~Del Duca, C.R.~Schmidt, Phys.\ Rev.\ D59 (1999) 074004, hep-ph/9810215.
\end{thebibliography}
\end{document}
